\section{Lambert resummation and nonlinear universality}
\label{sec:universality}

\subsection{A convergent chart for the scale itself}

The coefficient expansion is asymptotic. The scalar scale $B$ has a
stronger analytic structure that explains how it absorbs the
$\log\log n$ corrections.

\begin{proposition}[Convergent two-variable scale chart]
\label{prop:scale-chart}
Let $h=L^{-1}$, $w=L^{-1}\log L$, and
$c=\log(\ph(\lambda)/(2b))$. There is a holomorphic germ
$Y(h,w;\lambda,b)$, with $Y(0,0;\lambda,b)=1$, uniform on compact
real $(\lambda,b)$ sets as above, such that
\begin{equation}\label{eq:Y-scale}
 B^2=B_0^2Y(L^{-1},L^{-1}\log L;\lambda,b).
\end{equation}
It is determined by the removable holomorphic equation
\begin{equation}\label{eq:Y-equation}
\begin{split}
 Y={}&\frac{1-e^{-\lambda}E_2(2\lambda h)
                e^{-\lambda w-\lambda h(c+\log Y)}}{
                \lambda\ph(\lambda)}\\
 &+\frac{2h}{\ph(\lambda)}
   \left\{Z(2\lambda h)+\frac{\zeta(2+2\lambda h)^2}{b^2}\right\}.
\end{split}
\end{equation}
Its Taylor series in the two independent variables converges after
physical substitution for sufficiently large $L$.
\end{proposition}

\begin{proof}
The quotient is holomorphic at $\lambda=0$, since its numerator
vanishes there. At $h=w=0$, the right side equals
$(1-e^{-\lambda})/(\lambda\ph(\lambda))=1$ and its $Y$
derivative is zero. The implicit function theorem therefore applies
with derivative one after moving $Y$ to the left. A finite compactness
cover and uniqueness give a common neighborhood. Substituting
$B^2=B_0^2Y$ into \eqref{eq:scale} gives
\eqref{eq:Y-equation} exactly. The unique large positive solution
from \cref{thm:scale} identifies the branch.
\end{proof}

\subsection{Exact marginal Lambert scale}

At $\eps=0$, equation \eqref{eq:scale} is
\[
 \frac{bB^2}{n}=\log B+\frac32+\nu,
 \qquad \nu=\frac{\zeta(2)^2}{b^2}.
\]
Set $T=2bB^2/n$. Then
\begin{align}
 T-\log T&=L+3+2\nu-\log(2b),\label{eq:Lambert-scale}\\
 T&=-W_{-1}\left(-\frac{2b}{n}e^{-3-2\nu}\right),
 &B^2&=\frac n{2b}T,\qquad q=\frac2T.
       \label{eq:Lambert-explicit}
\end{align}
The lower real branch is required because $T\to\infty$
\cite{dlmf-lambert}. This is a scale for the coefficient law.
The coefficient scale and the earlier inverse-function Lambert core
have different constants; the following exact matching explains them.

\begin{corollary}[All marginal central coefficients as constants]
\label{cor:marginal}
Let $s=3$, $b=A=\zeta(2)$, and let $T>1$ solve
\begin{equation}\label{eq:critical-T}
 T-\log T=L+5-\log(2A).
\end{equation}
Define
\begin{equation}\label{eq:Dj}
 D_j=\frac1{\sqrt\pi\,j!}
   \Rea\left[(\partial_z+\log2+i\pi)^j\Gamma(z)
                                  \right]_{z=j+1/2},
 \qquad D_0=1.
\end{equation}
For each fixed $N$,
\begin{equation}\label{eq:marginal-T-series}
 R_b^n a_n=\sqrt{\frac A\pi}\,n^{-3/2}T^{-1/2}
       \left\{\sum_{j=0}^N\frac{D_j}{T^j}
                     +O(T^{-N-1})\right\}.
\end{equation}
If $Q_j(Y)$ denotes the Gamma-moment polynomial of
\cite{repo-moments}, then exactly
\begin{equation}\label{eq:matching}
 D_j=Q_j(-\log2).
\end{equation}
Thus the exact Lambert scale turns the previous logarithmic
polynomials into a sequence of constants.
\end{corollary}

\begin{proof}
At $\eps=0$ only $k=j$ survives in \eqref{eq:correction}.
Since $q=2/T$, the normalized central coefficient of $T^{-j}$ is
$(-1)^j\HH_{2j}^{[j]}(0)/(j!\normal(0))$.
Equation \eqref{eq:central-C} gives precisely \eqref{eq:Dj}.
The leading prefactor is
$\normal(0)/(nB)=\sqrt{A/\pi}\,n^{-3/2}T^{-1/2}$.
Finally the predecessor defines
\[
 Q_j(Y)=\frac1{\sqrt\pi\,j!}\Rea\left[
             (\partial_z-Y+i\pi)^j\Gamma(z)
                                  \right]_{z=j+1/2},
\]
so substitution gives \eqref{eq:matching}, with no asymptotic
comparison needed.
\end{proof}

For $C=\gamma+\log2$, the first constants are
\begin{align}
 D_1&=1-\frac C2,\notag\\
 D_2&=1-2C+\frac{3C^2}{8}-\frac{3\pi^2}{16}.
 \label{eq:D-first}
\end{align}
Re-expanding \eqref{eq:critical-T} in $L$ gives, in particular,
\begin{equation}\label{eq:marginal-old-first}
 R_b^n a_n=\sqrt{\frac A\pi}\,n^{-3/2}L^{-1/2}
 \left\{1-\frac{\log L-\log A+\gamma+3}{2L}
              +O\left(\frac{(\log L)^2}{L^2}\right)\right\}.
\end{equation}
This agrees with the integrated marginal result, including its
constant $\gamma+3$. At higher orders the same conclusion follows
from the exact identity \eqref{eq:matching} and the convergent scale
chart, rather than from a comparison of a few numerical values.

\subsection{Positive analytic outer functions}

The square root is not essential to the all-order logarithmic
calculus. Let $G$ vary in a family of holomorphic functions satisfying
the following uniform hypotheses for all sufficiently small real
$\eps$:
\begin{enumerate}[label=(\roman*)]
\item $G(0)=0$, and the Taylor coefficients of $G$ at zero are
nonnegative real numbers;
\item there are $\eta,M>0$ such that $G$ is holomorphic on
$|w|<\zeta(3+\eps)+\eta$ and bounded by $M$ there;
\item $G'(0)\ge\eta$.
\end{enumerate}
Shrinking $\eta$ makes it harmless to use the same symbol for the
uniform lower bound and the disk margin. Compact families with fixed
positive margins satisfy these conditions. Put
\begin{equation}\label{eq:G-parameters}
 f(z)=G(\Li_{3+\eps}(z)),\quad
 c=G'(\zeta(3+\eps)),\quad
 \kappa=\frac{A^2G''(\zeta(3+\eps))}{c},\quad \mu=cA.
\end{equation}
Then $c$ is bounded above and away from zero, and $\kappa$ is
nonnegative and bounded. Define $X_n$ using \eqref{eq:pgf} and set
\begin{equation}\label{eq:G-scale}
 \frac{B^2}{nc}=K_\eps(\log B)+\kappa=:V,
 \qquad q=\frac{E_2(\eps)B^{-\eps}}{V},\qquad
 \xi_m=\frac{m-n\mu}{B}.
\end{equation}

\begin{theorem}[Two-derivative composition principle]
\label{thm:universality}
Under the preceding uniform hypotheses, \cref{thm:all-orders}
holds with \eqref{eq:G-scale} in place of
\eqref{eq:scale}--\eqref{eq:q-xi}, and with precisely the same
functions $\mathcal C_j$. The leading scale is
\begin{equation}\label{eq:G-leading}
 B^2\sim\frac{ncL}{2}\ph(\lambda).
\end{equation}
For the germ $v(r)=f(re^{v(r)})$, write $a_n=[r^n]v(r)$ and
$R=e^{-f(1)}$. Then $nBR^n a_n$ has the same expansion at
$\xi=(n-n\mu)/B$.

After normalization by $f(1)$ and centering, all inverse-logarithmic
orders depend on the outer function only through its first two
derivatives at $\zeta(3+\eps)$. Higher derivatives contribute to the
remainder beyond every fixed inverse-logarithmic order.
\end{theorem}

\begin{proof}
The positive Taylor coefficients and the linear term of $G$ imply
$[z^k]f(z)\ge\eta k^{-3-\eps}$. Uniform boundedness on the larger
disk supplies uniform endpoint derivatives by Cauchy's estimates.
Writing $D=\zeta(s)-\Li_s(e^{-t})$, Taylor expansion at the endpoint
gives
\[
 G(\zeta(s)-D)-G(\zeta(s))
 =-cD+\frac12G''(\zeta(s))D^2+O(|D|^3).
\]
The proof of \cref{lem:local} therefore yields
\[
 f(e^{-t})-f(1)+\mu t
  =\frac{ct^2}{2}\{K_\eps(-\Log t)+\kappa\}
                                      +O(|t|^{5/2}).
\]
The scale existence proof applies with $b$ replaced by $1/c$ and
$\nu$ replaced by $\kappa$. The positive-intensity proof of
\cref{lem:fourier-tail} also applies verbatim with the indicated
uniform lower bound. The entire finite-expansion proof depends only
on these two facts, so it gives the same $\mathcal C_j$ and remainder.
Lagrange inversion proves the last coefficient identity exactly as in
\cref{prop:poisson}. Terms involving higher derivatives of $G$ begin
at $O(D^3)$ and are covered by the uniform local remainder.
\end{proof}

For the square-root family, choose the disk margin smaller than
$b_-^2/2$. Its branch point is at
$w=u^2/2=\zeta(s)+b^2/2$, and $G'(0)=1/u$ is uniformly positive.
Thus the main quadratic model satisfies the hypotheses, with uniform
boundedness on a slightly smaller common margin.

\begin{center}
\begin{tabular}{@{}lll@{}}
\toprule
Outer function $G(w)$ & $c$ & $\kappa$\\
\midrule
$a w$ & $a$ & $0$\\
$u-\sqrt{u^2-2w}$ & $b^{-1}$ & $A^2/b^2$\\
$a(e^{\beta w}-1)/\beta$ & $a e^{\beta\zeta(s)}$ & $\beta A^2$\\
$-(a/\beta)\log(1-\beta w)$ & $a/(1-\beta\zeta(s))$
 & $\beta A^2/(1-\beta\zeta(s))$\\
\bottomrule
\end{tabular}
\end{center}
Here $a>0$, $\beta>0$, and in the last row
$\beta\zeta(s)<1$ with a uniform margin. The limit $\beta=0$
recovers the linear row. Compact parameter subsets satisfying the
stated margins provide concrete examples of the theorem.

The principle does not say that these models have the same full
transseries. Their coefficients at exponential scales in $L$ can
depend on higher endpoint derivatives and on other singularities.
